\chapter*{Introduction}
\addcontentsline{toc}{chapter}{Introduction}
\markboth{Introduction}{Introduction}
\label{ch:introduction}

Excess flow valves (EFVs) are passive safety devices intended to limit or interrupt fluid discharge when the flow rate through a system exceeds a prescribed threshold. They are used in pressurized storage and distribution systems, where a downstream failure can otherwise lead to uncontrolled gas release. In a mechanically actuated valve, the pressure distribution and viscous stresses around the obturator produce a hydrodynamic load that must be balanced against the restoring force of the mechanical assembly. The relationship between the internal flow field and this load is therefore central to valve design.

Predicting the fluid forces is challenging because the flow can accelerate through restricted passages and develop strong velocity gradients, separation and recirculation. Experimental measurements can characterize the overall response but may provide limited access to the local three-dimensional flow structures. Computational Fluid Dynamics (CFD) offers a complementary approach for investigating pressure losses, local flow features and loads on the obturator \cite{white2016,moukalled2016}. Its predictions nevertheless depend on modeling assumptions, mesh resolution and numerical convergence, and should not be treated as experimentally validated without supporting evidence.

The objective of this thesis is to develop a repeatable CFD-based methodology for comparing four excess flow valve obturator configurations and investigating the influence of operating conditions. A full-factorial Design of Experiments (DoE) is defined by two gases (hydrogen, $\mathrm{H}_2$, and nitrogen, $\mathrm{N}_2$), three inlet velocity levels ($5$, $10$ and $15\,\mathrm{m\,s^{-1}}$) and two geometric parameters: number of holes (three or six) and hole diameter ($1.5$ or $3.0\,\mathrm{mm}$). The resulting design contains $2\times3\times2\times2=24$ prescribed combinations.

The principal research questions are: (i) how do obturator hole count and diameter affect the valve pressure loss and the axial hydrodynamic load; (ii) how do the working fluid and inlet velocity modify these responses; and (iii) how can the numerical information be organized into a reproducible basis for comparing candidate valve configurations? The responses of interest are the pressure difference between inlet and outlet, the force acting on the fixed obturator and their dimensionless counterparts.

The simulations use OpenFOAM 10 on the ARIES High-Performance Computing (HPC) infrastructure \cite{openfoam_user,hpc_aries}. The workflow combines geometry-specific mesh generation with scripted preparation, batch execution and post-processing. The longer-term goal is to make the resulting dataset accessible through an interactive comparison tool.

An important limitation is that the CFD model considers steady incompressible flow around a \emph{stationary} obturator. It can estimate hydrodynamic loads for prescribed configurations but does not directly simulate the transient closing motion, predict closure time or establish a validated trip threshold. Furthermore, local acceleration through restricted passages may challenge the incompressibility, so will be an important aspect to consider during the results analysis. 

The manuscript is organized as follows:
\begin{itemize}
    \item \textbf{Chapter 1} describes the governing equations, boundary conditions, mesh-generation strategy, solver configuration and numerical verification requirements.
    \item \textbf{Chapter 2} defines the factorial design, its geometric and operating factors, the automation workflow and the intended response-analysis methodology.
    \item \textbf{Chapter 3} presents the post-processing framework and will compare the CFD results at the different flow regimes.
    \item \textbf{Chapter 4} describes the proposed interactive application to compare valve configurations.
    \item \textbf{Chapter 5} will summarize the supported findings, limitations and possible future developments.
\end{itemize}
